% ---
% id: protocol-compiler/thm:feasibility
% kind: thm
% title: "Feasibility by an explicit or a verified candidate"
% summary: >-
%   An empty constraint list has an explicit product state in the feasible set.
%   Otherwise the compiler verifies, exactly and block by block, the default
%   candidate and, if it fails, a supplied witness, and refuses when neither
%   verifies.
% prominence: standard
% status: implemented
% uses: [def:feasible, def:compile, def:exact, def:homog, def:declaration]
% mentions: [thm:validator-soundness, thm:band-a, foundations/def:exact-island, req:feasibility-witness]
% assumes: []
% refs: []
% provenance: newly-derived
% implements: certiqs_key.protocol_compiler.feasibility:check_feasibility
% lean:
% notes: "The pointer to thm:band-a (i) at the one-block case is explanatory and is a mention. D(vi) is a hypothesis of (ii) alone: the constraint legs read each entry's relation symbol and right-hand side, which D(vi) types; (i) has no constraint leg."
% ---

D(vii) is the one target condition that is not a typing or a finite syntactic
check: whether $\Feas\ne\varnothing$ depends on the interaction of the marginal
with the constraint entries. It is therefore treated in two cases, and the
second one ends in a refusal rule rather than a theorem; the engine's order of
the two candidates and its refusal are Requirement~\ref{req:feasibility-witness}.

\begin{theorem}\label{thm:feasibility}
Let $\Decl$ be a candidate declaration satisfying D(i)--D(iii), and let
$\Feas$ be as in Definition~\ref{def:feasible}. (The hypothesis
is on the conditions alone: nothing below uses that $\Decl$ was emitted by
Definition~\ref{def:compile}, and Theorem~\ref{thm:validator-soundness} applies this theorem to
candidates.)
\begin{enumerate}[label=(\roman*),leftmargin=2.2em]
\item If $\CList=\varnothing$ then $\Feas\ne\varnothing$. Explicitly, the
\tref{block list}
\begin{equation}\label{eq:witness-empty}
  \omega^{\star(s)}=
  \begin{cases}
    \tau_A\otimes\dfrac{\idon{\Qsys_s}}{d_{\Qsys}},
      & \text{trusted-source mode},\\[1.2ex]
    \dfrac{\idon{A}}{d_A}\otimes\dfrac{\idon{\Qsys_s}}{d_{\Qsys}},
      & \text{untrusted-source mode},
  \end{cases}
  \qquad s\in\SecSet ,
\end{equation}
which denotes $\tau_A\otimes\idon{\Qsys}/d_{\Qsys}$, respectively
$\idon{A}/d_A\otimes\idon{\Qsys}/d_{\Qsys}$, lies in $\Feas$.
\item If $\CList\ne\varnothing$ and $\Decl$ satisfies D(vi) as well, which
makes each constraint leg of \eqref{eq:feas} a well-posed comparison, then
$\Feas\ne\varnothing$ holds if and only
if some block list $\omega$ satisfies the conditions of \eqref{eq:feas}. The
compiler decides those conditions \emph{exactly}, by finitely many arithmetic
operations and per block, on at most two candidates and in a fixed order:
first the \emph{default candidate} \eqref{eq:witness-empty}, which satisfies
every conjunct of \eqref{eq:feas} except possibly the constraint legs
(by the proof of (i)) and which the compiler forms in order to test and never
emits; and, if that candidate fails a constraint leg, a \emph{supplied}
witness $\omega^{\star}$, consumed and verified and never emitted. It refuses
exactly when the default candidate fails and either no witness is supplied or
the supplied one fails.
\end{enumerate}
\emph{The general format.} For a candidate declaration in the general format
of Definition~\ref{def:declaration}\,(DX1)--(DX6) satisfying D(i)--D(iii),
and D(vi) for part (ii), with $\Feas$ the set of \eqref{eq:feas-general},
both parts hold with the
\tref[joint sector]{joint sectors} in place of the sectors. In (i) the default candidate is
the block list over the joint sectors
\begin{equation}\label{eq:witness-general}
  \omega^{\star(t,s)}=
  \begin{cases}
    \tau^{(t)}_A\otimes\dfrac{\idon{\Qsys_s}}{d_{\Qsys}},
      & \text{trusted-source mode},\\[1.2ex]
    \dfrac{\idon{A_t}}{d_A}\otimes\dfrac{\idon{\Qsys_s}}{d_{\Qsys}},
      & \text{untrusted-source mode},
  \end{cases}
  \qquad (t,s)\in\mathbb{T}\times\SecSet ,
\end{equation}
and in (ii) the verification is the per-joint-block one listed in the proof,
which forms no array larger than one joint block. At $|\mathbb{T}|=1$ and one
\tref{measured factor} \eqref{eq:witness-general} is \eqref{eq:witness-empty}
and the two lists of decisions coincide item for item.
\end{theorem}

\begin{proof}
\emph{Idea.} The default candidate is $\tau_A\otimes\idon{\Qsys}/d_{\Qsys}$,
or $\idon{A}/d_A\otimes\idon{\Qsys}/d_{\Qsys}$ where there is no marginal,
presented block by block; every conjunct of the \tref{feasible set} except
the constraint legs holds of it by direct computation, and at an empty list
there are no constraint legs. It is formed and decided in
exact rational arithmetic and never emitted, so its entries need not be
binary64. Verifying any candidate is a finite list of exact decisions, each on
one block or on one $d_A\times d_A$ matrix. The general format repeats both
arguments over the joint sectors.

\emph{Step 1 (the default candidate lies in $\Feas$ at an empty list, (i)).}
Write $\chi_s\coloneqq\idon{\Qsys_s}/d_{\Qsys}$, so that
$\sum_s\tr\chi_s=\sum_sd_s/d_{\Qsys}=1$. The sector decomposition is not a
condition here: $\omega^{\star}$ is presented as a list, one array on
$A\Qsys_s$ per $s$, so the off-block conjunct holds by the presentation and
there is nothing to verify.

\emph{Positivity, per block.} In \tref{trusted-source mode}
$\omega^{\star(s)}=\tau_A\otimes\chi_s$ is Hermitian, $\tau_A$ being Hermitian
by D(ii) and $\chi_s$ a positive real multiple of an identity; and it is
positive semidefinite, being the tensor product of the positive semidefinite
$\tau_A$ (D(ii)) with the positive semidefinite $\chi_s$. In
\tref{untrusted-source mode} $\omega^{\star(s)}=\idon{A}/d_A\otimes\chi_s\succ0$ outright.
\emph{Trace.} $\tr\omega^{\star}=\sum_s\tr\omega^{\star(s)}
=\tr\tau_A\sum_s\tr\chi_s=\tr\tau_A=1$ in trusted-source mode, by D(ii); in
untrusted-source mode $\tr\omega^{\star}=1$ directly.
\emph{Marginal.} In trusted-source mode
$\omega^{\star}_A=\sum_s\tr_{\Qsys_s}\omega^{\star(s)}
=\tau_A\sum_s\tr\chi_s=\tau_A$; in untrusted-source mode there is no marginal
condition beyond the trace, already checked. \emph{Constraints.} With
$\CList=\varnothing$ the conjunction of \eqref{eq:feas-modes} is empty and
holds vacuously. Hence $\omega^{\star}\in\Feas$.

\emph{Step 2 (the number model of Step 1).} Two points of that argument
concern the number model. First, every step used D(ii) \emph{as an identity of
the declared numbers}, which is what Definition~\ref{def:exact} delivers ---
$\tr\tau_A=1$ with its sum rounding nowhere, in particular. Second,
\emph{the entries of $\omega^{\star}$ need not be \tref{binary64 numbers}, and
nothing requires them to be}: $1/d_{\Qsys}$ is a rational, dyadic only when
$d_{\Qsys}$ is a power of two, and $\omega^{\star}$ is \emph{never emitted}.
It is formed and decided in exact rational arithmetic on the declared dyadic
entries --- the exact arithmetic in which the \tref{exact island}
(\xref{foundations/def:exact-island}\,(i)) takes its decisions on
\tref{declared data}, applied here to a candidate the compiler forms itself
and never emits --- so every predicate above is decided, not approximated,
and no representability question arises for an
object no consumer ever sees. A \emph{supplied} witness is a different
matter: it is \tref{declared data}, hence dyadic, and its predicates are
decided exactly on those dyadic numbers.

\emph{Step 3 (the verification of a candidate, (ii)).} The stated equivalence
is Definition~\ref{def:feasible} read literally. The verification of a
candidate is finite and \emph{per block}. Writing $n_s\coloneqq d_Ad_s$ for
the size of the $s$-block, no array larger than one $n_s\times n_s$ block is
formed at any point; where $|\SecSet|=1$ that one block is $d_Ad_{\Qsys}$
square, and is a block and not an exception (Theorem~\ref{thm:band-a}\,(i)). It
consists of
\begin{itemize}[leftmargin=1.6em]
\item $|\SecSet|$ positivity decisions, each on an $n_s\times n_s$ matrix ---
  the Hermiticity decision $\omega^{(s)}=\omega^{(s)\dagger}$, entry by entry,
  and the semidefiniteness of the resulting Hermitian matrix by exact
  symmetric elimination, of which the largest is on $\max_sn_s$ and not on
  $d_Ad_{\Qsys}$;
\item one trace decision, $\tr\omega=1$ on the scalar
  $\sum_s\tr\omega^{(s)}$, a sum of $|\SecSet|$ block traces accumulated in the
  \tref{declared order};
\item in trusted-source mode one marginal decision: form
  $\omega_A=\sum_s\tr_{\Qsys_s}\omega^{(s)}$, which is
  $|\SecSet|$ partial traces onto $A$ and $|\SecSet|-1$ additions of $d_A^2$
  entries, then $\omega_A=\tau_A$ over $d_A^2$ entries;
\item and, per entry $k$, one constraint leg: in trusted-source mode
  $\sum_s\tr_{\Qsys_s}[(\idon{A}\otimes\hat O_k^{(s)})\omega^{(s)}]$, a
  matrix of size $d_A$ assembled from $|\SecSet|$ block terms, compared with
  $\Theta_k$ in the ordering $\bowtie_k$ --- an exact L\"owner comparison,
  i.e.\ the semidefiniteness of one $d_A\times d_A$ difference for an
  inequality, or one exact entrywise equality when $\bowtie_k$ is $=$; in
  untrusted-source mode the scalar
  $\sum_s\tr[(\idon{A}\otimes\hat O_k^{(s)})\omega^{(s)}]$ compared with
  $c_k$ as real numbers, according to the emitted source mode (D(vi)). The
  right-hand side of a trusted-source-mode leg is $c_k\tau_A$ for an entry
  tagged homogeneous and the declared $\Theta_k$ for one tagged general
  (Definition~\ref{def:homog}); the compiler forms $c_k\tau_A$ here in order
  to \emph{test} the candidate, and does not emit it --- the consuming engine
  constructs its own from the same declared $c_k$ and $\tau_A$, so the two
  form the same matrix.
\end{itemize}

\emph{Step 4 (one decision for an equality).} The equality decision is
equivalent to the two inequality conditions. For Hermitian operands $X,Y$, set
$D=X-Y$. If $D\succeq0$ and $-D\succeq0$, each real eigenvalue of $D$ is both
nonnegative and nonpositive, hence zero; Hermitian diagonalization gives
$D=0$. Conversely $D=0$ implies both inequalities. Thus entrywise $X=Y$
decides their conjunction exactly, in one predicate. The argument applies to
every operator row of the general format below as well.

All of these are finitely many arithmetic operations on the declared data, and
the largest object any of them forms is one $n_s\times n_s$ block or one
$d_A\times d_A$ matrix.

\emph{Step 5 (the general format, (i)).} Write
$\chi_s\coloneqq\idon{\Qsys_s}/d_{\Qsys}$ as above, so
$\sum_s\tr\chi_s=1$. The sector decomposition, now along the joint sectors,
is again the presentation and not a condition: \eqref{eq:witness-general} is
one array per joint sector. \emph{Positivity, per joint block}: in
trusted-source mode $\tau^{(t)}_A\otimes\chi_s$ is Hermitian and positive
semidefinite, each $\tau^{(t)}_A$ being so by D(ii) in its general-format
reading; in untrusted-source mode the block is positive definite outright.
\emph{Trace}: $\sum_{(t,s)}\tr\omega^{\star(t,s)}
=\big(\sum_t\tr\tau^{(t)}_A\big)\big(\sum_s\tr\chi_s\big)=1$ in
trusted-source mode by D(ii), and $\sum_t(d_{A_t}/d_A)\sum_s\tr\chi_s=1$ in
untrusted-source mode by $d_A=\sum_td_{A_t}$ (D(i)). \emph{Marginal}, per
Alice sector and in trusted-source mode:
$\sum_s\tr_{\Qsys_s}\omega^{\star(t,s)}
=\tau^{(t)}_A\sum_s\tr\chi_s=\tau^{(t)}_A$, which is the second line of
\eqref{eq:feas-general}; in untrusted-source mode there is no marginal
condition beyond the trace. \emph{Constraints}: with $\CList=\varnothing$ the
three constraint lines of \eqref{eq:feas-general} are empty conjunctions. So
$\omega^{\star}\in\Feas$, and the two points of Step 2 apply unchanged: every
step reads D(ii) as an identity of the declared numbers, and the entries of
\eqref{eq:witness-general} need not be binary64 because the candidate is
formed and decided in exact rational arithmetic and is never emitted.

\emph{Step 6 (the general format, (ii)).} The equivalence is
Definition~\ref{def:feasible} in its general-format reading, read literally.
Writing $n_{t,s}\coloneqq d_{A_t}d_s$ for the size of the joint block, the
verification of a candidate consists of
\begin{itemize}[leftmargin=1.6em]
\item $|\mathbb{T}|\,|\SecSet|$ positivity decisions, one Hermiticity and one
  exact symmetric elimination per joint block, the largest on
  $\max_{(t,s)}n_{t,s}$;
\item one trace decision on $\sum_{(t,s)}\tr\omega^{(t,s)}$, accumulated in
  the declared order of the joint sectors;
\item in trusted-source mode $|\mathbb{T}|$ marginal decisions, one per Alice
  sector: $|\SecSet|$ partial traces onto $A_t$, $|\SecSet|-1$ additions and
  one comparison, each over $d_{A_t}^2$ entries;
\item and, per entry $k$, its own leg. For an entry read through its
  \tref{lifted operator}, in trusted-source mode one exact L\"owner
  comparison per Alice sector on a $d_{A_t}\times d_{A_t}$ difference for an
  inequality, or one exact entrywise equality when $\bowtie_k$ is $=$, the
  right-hand side $c_k\tau^{(t)}_A$ for a homogeneous entry and the declared
  $\Theta^{(t)}_k$ for a general one; in untrusted-source mode one comparison
  of real numbers against $c_k$. For an entry declared party-local and scoped
  to $f$, one exact L\"owner comparison per pair $(t,s_{f^c})$ for an
  inequality, or one exact entrywise equality per pair when $\bowtie_k$ is
  $=$, on a difference of size
  $d_{A_t}\,d_{f^c,s_{f^c}}=n_{t,s}/d_{f,s_f}$, assembled from the
  $|\SecSet^{f}|$ joint blocks of that row by partial traces over
  $\Qsys_{f,s_f}$ --- which is the last line of \eqref{eq:feas-general}, and
  is \emph{never larger than one joint block}. The same decision rule applies
  in either source mode; a $1\times1$ row uses one real-number comparison as
  its scalar specialization.
\end{itemize}
\emph{No lifted operator is formed at any of these steps.} The value
$\tr_{\Qsys_s}[(\idon{A_t}\otimes\hat O^{\uparrow(s)}_k)\,\omega^{(t,s)}]$ is
computed by contracting the joint block against the declared block
$\hat O^{(s_{f_k})}_k$ on the tensor factor $\Qsys_{f_k,s_{f_k}}$ of that
block and tracing out the others, which reads one joint block and one block of
the scope and forms no array on $\Qsys$; the identity of that contraction with
the displayed trace is the definition of the lifted operator in (DX6).
\emph{So the largest object any decision of this verification forms is one
$n_{t,s}\times n_{t,s}$ joint block, one $d_{A_t}\times d_{A_t}$ matrix, or one
party-local difference, and each of the last two is at most one joint block};
at $|\mathbb{T}|=1$ and one measured factor the list is the list of Step 3,
item for item, with $n_{t_0,s}=n_s$.
\end{proof}
